% results.tex (v13, 2026-10-18 night: Fig. 4 c-e and Fig. 5 b-c added). Delta = RMSE(first) - RMSE(second) on the same site-years; 95% paired block-bootstrap interval in brackets.
\section{Results}\label{sec:results}

\subsection{Effect of the year-specific climate input}
The observations are unevenly distributed in time (Fig.~\ref{fig:data}b, Table~\ref{tab:data}). In Alaska, \ntopyrshareAk\% of the site-years were observed in \ntopyrAk, whereas the Lena Delta labels are spread over \nnyrLe{} years. Most measurements were made in August (median day of year \ndoymedAk{} in Alaska, \ndoymedLe{} in the Lena Delta and \ndoymedCa{} in Canada), before the end of the thaw season, and the TDD accumulated to the observation day amounted to a median of \ntddfracAk, \ntddfracLe{} and \ntddfracCa{} of the full-season sum (equation~(\ref{eq:tddobs})).

Replacing the mean climate input by the year-specific input reduced the spatial-block RMSE of the recalibrated Stefan model in the Lena Delta ($\Delta$RMSE \ndforcLePo{} cm) and that of the physics-residual model in the Lena Delta (\ndforcLeRo{} cm) and Canada (\ndforcCaRo{} cm) (Fig.~\ref{fig:accuracy}a). In Alaska the change was not resolved (Stefan \ndforcAkPo{} cm), consistent with the concentration of the labels in one year, for which the year-specific input adds little information beyond the spatial pattern. The year-specific input also ties each prediction to the climate of a stated year, which is the prerequisite for the yearly maps presented below.

\paperfig{Fig3_accuracy}{1.0}{\textbf{Year-specific climate input, validation design and conditions of the residual gain.} All panels use the site-year labels of 2010--2024 and out-of-fold predictions; the year and block-by-year designs score only years with at least 20 site-years. \textbf{a}, Spatial-block RMSE with the mean climate input (2011--2024 mean) and with the year-specific climate input (observation year, accumulated to the observation day). \textbf{b}, RMSE under five validation designs; under region hold-out the recalibrated Stefan model equals the source-coefficient model, and the value beyond the axis is printed. \textbf{c}, $\Delta$RMSE between the physics-residual model and the recalibrated Stefan model for six covariate sets (runs listed in Supplementary Table~S4). \textbf{d}, The same $\Delta$RMSE within regional terciles of eight covariates; outlined cells have 95\% intervals excluding zero, and hatched cells contain fewer than 30 site-years or five blocks. Intervals are 95\% paired block-bootstrap intervals (2000 resamples of 0.5$^\circ$ blocks); negative values favour the physics-residual model.}{fig:accuracy}

\subsection{Accuracy under spatial, temporal and regional validation}
Model rankings depended on the validation design (Fig.~\ref{fig:accuracy}b, Supplementary Table~S5). Random cross-validation yielded the lowest errors for direct ML (\nrmseAkrandDz{} cm in Alaska), and spatial-block cross-validation increased the direct-ML error by \noptimAkDz, \noptimLeDz{} and \noptimCaDz{} cm in Alaska, the Lena Delta and Canada. Year hold-out increased the direct-ML error in the Lena Delta from \nrmseLeblkDz{} to \nrmseLeyrDz{} cm, reflecting the spread of its labels over many years, and the block-by-year design produced the largest within-region errors for all methods fitted to labels of the region in all three regions. When an entire region was withheld, the direct-ML error in the Lena Delta rose to \nrmseLeregDz{} cm, compared with \nrmseLeregPz{} cm for the Stefan model with the coefficient of the other regions ($\Delta$RMSE \ndrefLeregDz{} cm). In the same test the physics-residual model remained close to the Stefan model (\ndrefLeregRo{} cm) and improved on it in Canada (\ndrefCaregRo{} cm). Spatial-block cross-validation is therefore used as the main design and region hold-out as the test of transferability.

\subsection{Gain of physics-residual learning and its controls}
Under spatial-block cross-validation with the year-specific input, the physics-residual model reduced the RMSE of the recalibrated Stefan model from \nrmseLeblkPo{} to \nrmseLeblkRo{} cm in the Lena Delta ($\Delta$RMSE \ndrefLeblkRo{} cm) and from \nrmseCablkPo{} to \nrmseCablkRo{} cm in Canada (\ndrefCablkRo{} cm), whereas in Alaska the two models reached \nrmseAkblkPo{} and \nrmseAkblkRo{} cm (\ndrefAkblkRo{} cm). The gain depended on the covariates available to the residual learner (Fig.~\ref{fig:accuracy}c). With climate covariates only, the residual model did not improve on the Stefan model in Alaska (\ndfeatAkFclim{} cm) or Canada (\ndfeatCaFclim{} cm); adding terrain and soil covariates produced the gains above. MODIS vegetation and snow timing (2010--2024) or high-resolution soil and land-cover classes changed the result by less than 0.5 cm in the Lena Delta and Canada, and the latter added a small gain in Alaska (\ndfeatAkFhires{} cm).

The result depended on how the physical model entered the learner (Supplementary Fig.~S2). Augmenting the training set with Stefan pseudo-labels reproduced the Stefan model (Lena Delta \ndrefLeblkDo{} cm), and using the Stefan prediction as an input feature yielded results similar to direct ML. With all labels of a region, direct ML reached lower point errors than the physics-residual model (Alaska \nrmseAkblkDz{} versus \nrmseAkblkRo, Lena Delta \nrmseLeblkDz{} versus \nrmseLeblkRo{} and Canada \nrmseCablkDz{} versus \nrmseCablkRo{} cm), but its difference from the recalibrated Stefan model was resolved only in Canada (\ndrefCablkDz{} cm), and its error increased under region hold-out. Replacing CatBoost by random forest, a multilayer perceptron, TabM or FT-Transformer did not change the sign of the physics-residual gain in the Lena Delta and Canada.

Across regional terciles of eight covariates, the gain occurred in most conditions of the Lena Delta (\ngcLesig{} of \ngcLecells{} terciles with intervals excluding zero) and Canada (\ngcCasig{} of \ngcCacells) and increased with the error of the Stefan model; across terciles, the correlation between the Stefan error and the gain was \ngcLecorr{} in the Lena Delta and \ngcCacorr{} in Canada (Fig.~\ref{fig:accuracy}d). In the warmest TDD tercile of the Lena Delta the gain reached \ndtddLewarm{} cm, whereas no tercile in Alaska changed by more than \ngcAkmaxabs{} cm. These patterns indicate that the residual learner corrects departures of the effective Stefan coefficient from its regional value that arise from soil and terrain. Consistent with this interpretation, terrain variables accounted for \nshapLeResTerrain\% of the mean absolute SHAP value of the residual model in the Lena Delta and soil variables for \nshapCaResSoil\% in Canada (Supplementary Fig.~S3).

\subsection{Sample efficiency and transfer to new regions}
The label-number experiments show how the methods behave when labels are scarce (Fig.~\ref{fig:labels}a). With 25 training labels drawn from the 2010--2024 site-years of a region, the physics-residual model was at least as accurate as a Stefan model recalibrated on the same labels (Alaska \nllAkRotwfive, Lena Delta \nllLeRotwfive{} and Canada \nllCaRotwfive{} cm), and with 100 labels its gain was resolved in all three regions (\nllAkRohundred, \nllLeRohundred{} and \nllCaRohundred{} cm). Direct ML trained on 25 labels had a higher point error than the Stefan model in Alaska (\nllAkDztwfive{} cm) and Canada (\nllCaDztwfive{} cm). Physics-input ML and physics pseudo-labels followed the patterns of their full-data results (Supplementary Fig.~S4).

\paperfig{Fig4_labels}{1.0}{\textbf{Accuracy with few labels, in new regions and in external CALM regions.} \textbf{a}, $\Delta$RMSE between each method and a Stefan model recalibrated on the same $n$ training labels within each region (labels of 2010--2024, three draws whose predictions are averaged before scoring; $n$ extends to the labels available in each region); the band is the 95\% block-bootstrap interval of the physics-residual model. \textbf{b}, RMSE in each region for models trained on the other two regions plus $n$ labels from the target region; a value beyond the axis is drawn at the top edge and printed. \textbf{c}, External CALM sites, numbered by sub-region, and the label cells of the three study regions. \textbf{d}, RMSE in the external regions (observations of 2010--2020) against the number of labels from the same external region, on the \nftrows{} site-years with values of Wei et al., Liu et al. and ESA CCI v5 (2010--2020); horizontal lines mark the products, of which Wei et al. and Liu et al. include these sites in their training data; the pooled set covers the seven sub-regions with product values. \textbf{e}, RMSE of the physics-residual model with no and with five local labels, and of ESA CCI v5, in the sub-regions with at least 30 such site-years (numbers as in \textbf{c}; site-years in brackets).}{fig:labels}

In the transfer experiment, each region was predicted by models trained on the other two regions and $n$ labels from the target region (Fig.~\ref{fig:labels}b). Without target labels, direct ML reached \ntlrmseLeDszero{} cm in the Lena Delta, compared with \ntlrmseLePzzero{} cm for the source-coefficient Stefan model and \ntlrmseLeRszero{} cm for the physics-residual model. Twenty-five target labels reduced the error of the physics-residual model to \ntlrmseLeRstwfive{} cm and that of direct ML to \ntlrmseLeDstwfive{} cm. In Alaska the physics-residual model gave the lowest error at every $n$ (\ntlrmseAkRszero{} cm without target labels). In Canada, direct ML trained on the source regions was the most accurate (\ntlrmseCaDszero{} cm); recalibrating the Stefan coefficient on 25 Canadian labels increased the error from \ntlrmseCaPzzero{} to \ntlrmseCaPntwfive{} cm, and the physics-residual model, which uses this anchor, followed (\ntlrmseCaRszero{} to \ntlrmseCaRstwfive{} cm), reflecting the wide range of thaw depths among the Canadian sites.

The same procedure was applied to eight external CALM regions with \nftScoredRows{} scored site-years from 2010--2024 that were not used in model development (Fig.~\ref{fig:labels}c--e, Supplementary Fig.~S5). On the \nftrows{} site-years of seven regions with values of Wei et al., Liu et al. and ESA CCI v5, the physics-residual model trained on the three study regions reached \nftAllRszero{} cm without local labels. Five labels from the same external region reduced the error to \nftAllRsfive{} cm and 25 labels to \nftAllRstwfive{} cm, whereas direct ML reached \nftAllDstwfive{} cm with 25 labels. With five labels the physics-residual model was more accurate than ESA CCI v5, which is not trained on ALT observations (\nftAllCcifivezero{} cm on the same \nftrows{} site-years; $\Delta$RMSE \nftdRsCcifivefive{} cm on all \nftdCciRows{} site-years with an ESA CCI v5 value). Wei et al. and Liu et al., which include these sites in their training data, reached \nftAllWeizero{} and \nftAllLiuzero{} cm. Most of the reduction came from recalibrating the Stefan coefficient; the residual learner did not reduce the error further at these sites (\nftdRsPnfive{} cm with five labels on all \nftScoredRows{} site-years), where the active layers of West and Central Siberia and Svalbard exceed the range of the study regions. The pooled values combine different regional responses (Fig.~\ref{fig:labels}e): with five labels the error decreased in West Siberia, Central Siberia, European Russia and Svalbard, but increased in Chukotka (\nftChRszero{} to \nftChRsfive{} cm) and Kamchatka (\nftKaRszero{} to \nftKaRsfive{} cm; 16 site-years, not shown in Fig.~\ref{fig:labels}e), where the error without local labels was below 10 cm, and in north-eastern Siberia (\nftNsRszero{} to \nftNsRsfive{} cm).

\subsection{Comparison with published products in the same cells and years}
On the common site-years (\nyrscc), the physics-residual model reached \nccLeRoall{} cm in the Lena Delta, compared with \nccLeCcifiveall{} cm for ESA CCI v5, \nccLeAaltoall{} cm for Aalto et al., \nccLeCcifourall{} cm for ESA CCI v4, \nccLeLiuall{} cm for Liu et al. and \nccLeWeiall{} cm for Wei et al. (Fig.~\ref{fig:products}a, Supplementary Table~S6). On the Lena Delta site-years more than 25 km from the training sites of Wei et al. (\nnprodLeweimask{} site-years with a value of Wei et al., \nnprodLeccifivemask{} with a value of ESA CCI v5), the paired differences were \ndprodLeweiRomask{} cm relative to Wei et al. and \ndprodLeccifiveRomask{} cm relative to ESA CCI v5. In Alaska (\nccAkRoall{} versus \nccAkWeiall{} and \nccAkLiuall{} cm for Wei et al. and Liu et al.) and Canada (\nccCaRoall{} versus \nccCaWeiall{} and \nccCaLiuall{} cm), the model and the two ML products performed on par.

On the Lena Delta site-years, the two ML products assign most cells values between about 40 and 70 cm regardless of the observed ALT (Fig.~\ref{fig:products}b): the correlation between observed and predicted values was \nrLeWei{} for Wei et al. and \nrLeLiu{} for Liu et al., against \nrLeRo{} for the physics-residual model, and the mean bias was \nbiasLeAllWei{} and \nbiasLeAllLiu{} cm, against \nbiasLeAllRo{} cm. The bias of all series increased towards thin active layers (Fig.~\ref{fig:products}c): for cells with an observed ALT below 30 cm (\nthinNLe{} site-years), Wei et al. and Liu et al. overestimated by \nbiasLeThinWei{} and \nbiasLeThinLiu{} cm and the physics-residual model by \nbiasLeThinRo{} cm. In Alaska the same class (\nthinNAk{} site-years) was overestimated by \nbiasAkThinWei{}, \nbiasAkThinLiu{} and \nbiasAkThinRo{} cm, and the bias of the physics-residual model, Wei et al. and Liu et al. decreased with the observed ALT.

\paperfig{Fig5_products}{1.0}{\textbf{Comparison with published products in the same cells and years.} \textbf{a}, RMSE of three methods (spatial-block out-of-fold predictions) and five products on the common site-years (\nyrscc), for all cells and for cells more than 25 km from the training sites of Wei et al.; ESA CCI v5 and Aalto et al. are scored where they have a value, and values above 45 cm are drawn on the compressed right segment of the axis. \textbf{b}, Observed against predicted ALT on the \nccLeN{} common site-years of the Lena Delta for the physics-residual model (spatial-block out-of-fold) and the two ML products, with the 1:1 line and the Pearson correlation; RMSE and mean bias are given in the text. \textbf{c}, Mean bias within classes of observed ALT on the common site-years of Alaska and the Lena Delta; numbers below the axis give the site-years per class, classes with fewer than 10 site-years are not drawn, ESA CCI v5 is scored where it has a value, and a value beyond the axis is drawn at the edge and printed. \textbf{d}, RMSE per observation year on the site-years of each year with values of the physics-residual model, the recalibrated Stefan model and Wei et al. (a larger set than in \textbf{a}--\textbf{c}), for years with at least 30 site-years (Alaska \nyrsywAk; Lena Delta \nyrsywLe; Canada \nyrsywCa); Liu et al. end in 2020 and are omitted in years with fewer than 30 values; bars give the number of site-years per year. The key below \textbf{c} applies to \textbf{b}--\textbf{d}.}{fig:products}

The regional contrast is consistent with the training data of the products. Alaska label cells lie at a median distance of \ndistweiAkmed{} km from the training sites of Wei et al., and Liu et al. were trained on an earlier version of the same Alaska and Canada field data, whereas the Lena Delta cells lie at a median distance of \ndistweiLemed{} km (Supplementary Fig.~S6). Both products overestimate the thin active layers of the northern delta, with mean biases of \nbiasLeNorthWei{} and \nbiasLeNorthLiu{} cm north of 72$^\circ$N. A model trained without any Lena Delta label was already more accurate there than Wei et al. on the same site-years (\ntlvprmseLeRsWeizero{} versus \ntlvprmseLeWeizero{} cm; $\Delta$RMSE \ntlvpLeRsWeizero{} cm), indicating that the advantage in this region does not depend on local labels alone.

Maps of one common year show these differences in space (Fig.~\ref{fig:prodmaps}); 2013 is the only year with observations in all three regions. In Alaska, the three maps have similar errors at the \nfeAkN{} cells observed in 2013 with values in all three maps (RMSE \nfeAkOof{} cm for the out-of-fold predictions of our model, \nfeAkWei{} cm for Wei et al. and \nfeAkLiu{} cm for Liu et al.); \nfeAkNearWei{} of these cells lie within 25 km of a training site of Wei et al., so the product values there are not independent of the observations. The largest differences between the maps occur in mountain terrain above 600 m, which lies largely outside the area of applicability and where the products are deeper than our map by \nmapAkFootMtnWei{} and \nmapAkFootMtnLiu{} cm between 68 and 70$^\circ$N. In the Lena Delta, both products are deeper than the observations across the delta (mean bias \nfeLeBiasWei{} and \nfeLeBiasLiu{} cm at the \nfeLeN{} cells observed in 2013; RMSE \nfeLeWei{} and \nfeLeLiu{} cm), whereas our model reaches an RMSE of \nfeLeOof{} cm; the number of cells is small in this year, and the site-year comparison of Fig.~\ref{fig:products} covers all years. In Canada (NW), the RMSE at the \nfeCaN{} observed cells was \nfeCaOof{} cm for our model, \nfeCaWei{} cm for Wei et al. (mean bias \nfeCaBiasWei{} cm) and \nfeCaLiu{} cm for Liu et al., whose training data include an earlier version of these Canada observations.

\paperfig{Fig6_product_maps}{1.0}{\textbf{ALT maps of the physics-residual model and two published products for 2013, the only year with observations in all three regions.} \textbf{a--c}, Alaska: physics-residual model, Wei et al. and Liu et al. on one colour range for all maps of the figure (1st--99th percentiles of the pooled values of our three maps; values beyond it are drawn in the end colours), with only our map faded outside its area of applicability. \textbf{d}, Predicted minus observed ALT at the \nfeAkN{} cells observed in 2013 with values in all three maps, and the RMSE of each map; for our model the predictions are spatial-block out-of-fold values with the TDD to the observation day, whereas the maps use the TDD to 1 September and the model fitted to all labels. Values beyond the axis are printed. Product values are not masked for training sites; \nfeAkNearWei{} of the \nfeAkN{} cells lie within 25 km of a training site of Wei et al. \textbf{e--h}, The same for the Lena Delta (\nfeLeN{} cells, \nfeLeFarWei{} of them more than 25 km from these training sites). \textbf{i--l}, The same for Canada (NW) (\nfeCaN{} cells, \nfeCaNearWei{} within 25 km). The map domain is the area with an ESA CCI permafrost probability of at least 50\%; the 1 km values are interpolated to 250 m (Lena Delta) or 500 m (Alaska, Canada) for display.}{fig:prodmaps}

Year by year, the physics-residual model had a lower error than both Wei et al. and the Stefan model in all \nywLen{} Lena Delta years with at least 30 site-years (\nyrsywLe), in \nywAkRoWeiwin{} of \nywAkRoWein{} Alaska years against Wei et al. and \nywAkRoPowin{} of \nywAkRoPon{} against the Stefan model (\nyrsywAk), and in \nywCaRoPowin{} of \nywCaRoPon{} Canadian years against the Stefan model (Fig.~\ref{fig:products}d). At Alaska cells observed in several years, the correlation between observed and predicted year-to-year changes was \nrdeltaAkRo{} for the physics-residual model and \nrdeltaAkWei{} and \nrdeltaAkLiu{} for Wei et al. and Liu et al., whereas in the Lena Delta the products followed the interannual changes more closely (\nrdeltaLeWei{} and \nrdeltaLeLiu{} versus \nrdeltaLeRo).

The residual learner can also be applied with a product as the anchor (equation~(\ref{eq:prod})). With the Lena Delta labels, adding the residual learner to the recalibrated product of Wei et al. changed its error by \npadLeRpPcal{} cm, and the Stefan-anchored model remained more accurate than the product-anchored model (\npadLeRoRp{} cm; Supplementary Fig.~S7). In Alaska and Canada the two anchors were not distinguishable.

\subsection{Year-labelled 1 km ALT maps, 2011--2024}
The physics-residual models fitted to all site-years of each region were applied with the ERA5-Land climate of each year from 2011 to 2024, with the TDD accumulated to 1 September (Fig.~\ref{fig:maps}). The maps cover \nncellmapAk{} cells in Alaska, \nncellmapLe{} in the Lena Delta and \nncellmapCa{} in Canada (NW). Relative to the 2011--2024 mean, the regional mean anomaly ranged from \nanminAk{} cm in \nanminyrAk{} to \nanmaxAk{} cm in \nanmaxyrAk{} in Alaska, from \nanminLe{} cm in \nanminyrLe{} to \nanmaxLe{} cm in \nanmaxyrLe{} in the Lena Delta and from \nanminCa{} cm in \nanminyrCa{} to \nanmaxCa{} cm in \nanmaxyrCa{} in Canada, with spatial patterns that differ between years (Fig.~\ref{fig:maps}f). In the Lena Delta, the mean anomaly of the observed cells had the same sign as the predicted regional anomaly in \nanobssameLe{} of \nanobsnLe{} years with at least 50 site-years (\nanobsyrsLe; $r$ = \nanobsrLe), and the correlation remained \nanobsrLeNoTwenty{} without 2020, when the observed anomaly of \nanobsLeTwentyN{} site-years reached \nanobsLeTwentyObs{} cm. In Alaska, where the set of observed cells changes between campaigns, the agreement was weaker ($r$ = \nanobsrAk), and the predicted interannual range (\nanpredrangeAk{} cm) was smaller than that of the observed cell means (\nanobsrangeAk{} cm); in Canada the correlation over \nanobsnCa{} years was \nanobsrCa, with the observed anomalies of 2023 and 2024 (80 site-years each) far above the predicted range (Fig.~\ref{fig:maps}d). The linear trend of the yearly maps over 2011--2024 averaged \ntrendAk{}, \ntrendLe{} and \ntrendCa{} cm per decade over the Alaska, Lena Delta and Canada cells, with positive values in \ntrendposAk\%, \ntrendposLe\% and \ntrendposCa\% of the cells (Fig.~\ref{fig:maps}e); the trend reflects the ERA5-Land climate of the period as transmitted by the model and is not an observed trend. The area of applicability covers \naoaAk\% of the Alaska cells, \naoaLe\% of the Lena Delta cells and \naoaCa\% of the Canada cells; cells outside it lie mainly in mountains and at the southern margin of the permafrost zone (Supplementary Fig.~S8). The half-width of the 90\% prediction interval is \nhwAk{}, \nhwLe{} and \nhwCa{} cm. The correction added by the residual learner is mostly within $\pm$10 cm (Supplementary Fig.~S9), and series at long-record cells are shown in Supplementary Fig.~S10.

\paperfig{Fig7_maps}{1.0}{\textbf{Year-labelled 1 km ALT maps, 2011--2024.} \textbf{a}, Mean ALT of 2011--2024 from the physics-residual model in Alaska, computed from the yearly maps; each yearly map uses the ERA5-Land climate of that year with the TDD accumulated to 1 September. Lighter colours mark cells outside the area of applicability and grey land lies outside the map domain. Two 50 km windows (North Slope near Toolik; Interior near Fairbanks) are enlarged at the right with the 1 km cells drawn without interpolation. \textbf{b}, The same for the Lena Delta, with a 24 km window around the Samoylov and Kurungnakh sites. \textbf{c}, The same for Canada (NW), with 50 km windows at Old Crow Flats and in the Mackenzie Delta near Inuvik; a--c share one colour range. \textbf{d}, Regional mean anomaly of the predicted maps and mean anomaly of the observed cells in years with at least 50 site-years (Alaska \nanobsyrsAk, Lena Delta \nanobsyrsLe, Canada \nanobsyrsCa); observed anomalies are taken relative to the 2011--2024 map mean of each observed cell and centred on the mean predicted anomaly of the same years, and the Lena Delta 2020 value is labelled with its number of site-years; the y ranges differ between regions. \textbf{e}, Linear trend of the yearly maps over 2011--2024 (least squares per cell, cm per decade) in the three regions; the trend describes the modelled response to the ERA5-Land climate of each year and is not an observed trend. \textbf{f}, ALT anomaly of each year relative to the 2011--2024 mean in Alaska, on a common colour range ($\pm$15 cm). The maps of a--c are interpolated from 1 km to 250 m for display.}{fig:maps}
